\begin{table}[H]
\centering
\caption{Standardized Effect Sizes for Main Outcomes}
\label{tab:sde}
\begin{adjustbox}{max width=\textwidth}
\small
\begin{tabular}{llccccccc}
\toprule
Outcome & Specification & $\hat{\beta}$ & SD($X$) & SD($Y$) & SDE & SE(SDE) & Classification \\
\midrule
Injury rate & Mine + Quarter FE & -0.0367 & 9.47 & 33.58 & -0.0103 & 0.0041 & Small negative \\
Serious inj.\ rate & Mine + Quarter FE & -0.0266 & 9.47 & 23.59 & -0.0107 & 0.0055 & Small negative \\
\bottomrule
\end{tabular}
\end{adjustbox}
\par\vspace{0.3em}
{\footnotesize \emph{Notes:} This table reports standardized effect sizes (SDE) to facilitate cross-study comparison of treatment effect magnitudes.
For continuous treatments, SDE $= \hat{\beta} \times \text{SD}(X) / \text{SD}(Y)$, which gives the effect of a one-standard-deviation change in the treatment variable, measured in standard deviations of the outcome.
SD($Y$) and SD($X$) are unconditional standard deviations from the summary statistics (Table~\ref{tab:summary}), before conditioning on fixed effects.

\textbf{Country:} United States.
\textbf{Research question:} Whether MSHA's 2007 civil penalty reform (30 CFR Part 100), which raised average proposed penalties 4.2-fold overnight, reduced mine-level injury rates.
\textbf{Policy mechanism:} The reform increased penalty assessment points across the board for all mine safety violations, with disproportionate increases for Significant and Substantial (S\&S) violations; it raised the financial cost of non-compliance to deter unsafe practices.
\textbf{Outcome definition:} Quarterly injury count per mine divided by mine employment, multiplied by 100, using MSHA accident/injury reports.
\textbf{Treatment:} Continuous; mean proposed penalty per S\&S violation at each mine during 2004--2006, divided by 100.
\textbf{Data:} MSHA Accidents, Violations, and Mines datasets, 2004--2010, mine-quarter panel.
\textbf{Method:} Continuous-treatment DiD with mine and quarter fixed effects, standard errors clustered at the mine level. N = 129,220 mine-quarter observations across 4,615 mines.
\textbf{Sample:} Active mines with at least one S\&S violation during 2004--2006 and non-missing employee counts.

Classification thresholds (7 categories): large negative ($< -0.15$), moderate negative ($-0.15$ to $-0.05$), small negative ($-0.05$ to $-0.005$), null ($-0.005$ to $0.005$), small positive ($0.005$ to $0.05$), moderate positive ($0.05$ to $0.15$), large positive ($> 0.15$).
Classification is based solely on the SDE point estimate --- never on statistical significance or p-values.
Classification labels refer to the magnitude of the standardized point estimate, not to statistical significance. ``Null'' denotes a near-zero effect size ($|$SDE$| < 0.005$), not a failure to reject a null hypothesis.}
\end{table}
